\documentclass{article}
\usepackage{amsmath,amsthm}
\newtheorem{lemma}{Lemma}
\newtheorem{theorem}{Theorem}

\title{Geometric response profiles near a limiting boundary}
\author{}
\date{}

\begin{document}
\maketitle

\section{Convention and stability criterion}

Throughout, the observation window is
\[
 \mathcal W=\{x\in\mathbb R:0\leq x<1\}.
\]
A sequence of nonnegative profiles will be called \emph{stable} when the
following uniform requirement holds: given any positive tolerance, one can
choose a single sufficiently late stage, depending on the tolerance but not
on the observation point, after which every later profile stays below that
tolerance at every point of the observation window.

For \(n\geq1\) put \(b_n(x)=x^n\).

\section{Local decay}

\begin{lemma}\label{lem:local-decay}
At every observation point \(x\in\mathcal W\), the values \(b_n(x)\) tend
to zero.
\end{lemma}

\begin{proof}
If \(x=0\) this is immediate. Otherwise \(0<x<1\), so
\(n\log x\to-\infty\), and hence \(x^n=\exp(n\log x)\to0\).
\end{proof}

\begin{theorem}\label{thm:uniform-decay}
The geometric profiles \(b_n\) are stable.
\end{theorem}

\begin{proof}
Fix a positive tolerance. Lemma~\ref{lem:local-decay} supplies a sufficiently
late stage at every observation point, and continuity preserves each such
estimate on a neighbourhood of that point. The resulting neighbourhoods
cover the bounded observation window, so retain finitely many and take the
largest associated stage. Every observation point then lies in one retained
neighbourhood, and all later profiles are smaller still. This is precisely
stability.
\end{proof}

\end{document}
